\section{Security Analysis}


\subsection{Threat Model}

We consider two adversary classes:

\begin{enumerate}
\item \textbf{Classical adversary} $\mathcal{A}_C$: Polynomial-time, controls $f < t$ validators.
\item \textbf{Quantum adversary} $\mathcal{A}_Q$: Has access to a fault-tolerant quantum computer, controls $f < t$ validators.
\end{enumerate}

\begin{theorem}[Key Extraction Resistance]
Under the MLWE assumption with parameters $(n, k, q, \eta)$ as specified in FIPS 203/204, neither $\mathcal{A}_C$ nor $\mathcal{A}_Q$ can extract the complete private key from fewer than $t$ shares.
\end{theorem}

\begin{proof}
Reconstructing the secret from $f < t$ shares requires solving $f$ equations in $t$ unknowns over $R_q^k$---information-theoretically impossible by Shamir's threshold property.
\end{proof}

Proactive share refresh (Section~5.3) provides forward secrecy: compromising shares in epoch $e$ does not help after refresh in $e+1$. During the BLS/ML-DSA migration period, hybrid signatures provide defense in depth---security holds if \emph{either} the discrete log assumption (BLS) \emph{or} the MLWE/MSIS assumption (ML-DSA) holds.

